\documentclass[11pt]{article}
\usepackage[margin=0.9in]{geometry}
\usepackage[T1]{fontenc}
\usepackage{newtxtext,newtxmath}
\usepackage{microtype}
\usepackage{amsmath,bm}
\usepackage{booktabs,tabularx,array}
\usepackage{enumitem}
\usepackage{graphicx}
\usepackage{tikz}
\usetikzlibrary{arrows.meta,positioning}
\usepackage{natbib}
\usepackage{hyperref}
\usepackage{xurl}
\hypersetup{hidelinks,pdfauthor={Rajeev Yadav, Ph.D.},pdftitle={FinCompass v2.0: Separating Bayesian Model Validity from Predictive Evidence in Probabilistic Market Forecasting}}
\setlength{\parindent}{1.1em}
\setlength{\parskip}{0pt}
\setlength{\emergencystretch}{2em}
\newcommand{\Prb}{\mathbb{P}}
\newcommand{\I}{\mathbb{I}}
\newcommand{\sigmoid}{\operatorname{sigmoid}}
\newcommand{\logit}{\operatorname{logit}}
\title{\textbf{FinCompass v2.0: Separating Bayesian Model Validity from Predictive Evidence in Probabilistic Market Forecasting}}
\author{Rajeev Yadav, Ph.D.\\\texttt{rajeevyadav@gmail.com}}
\date{August 2026}
\begin{document}
\maketitle
\begin{center}
\textbf{Version 2.0}
\end{center}

\begin{abstract}
Financial forecasting systems often collapse two different questions into one pass/fail decision: whether a probabilistic model is mathematically and operationally valid, and whether it has demonstrated strong out-of-sample predictive skill. This paper develops a layered architecture that separates those questions. A regularized Bayesian logistic reference model can be retained as a \emph{Bayesian baseline} when it satisfies hard validity conditions---chronological construction, matured targets, finite posterior quantities, coherent probabilities, adequate estimation support, and explicit applicability---even if it does not satisfy stronger locked-test skill gates. Stronger tiers, \emph{validated research} and \emph{validated market}, preserve stricter requirements based on proper scoring rules, calibration, temporal coverage, bootstrap uncertainty, walk-forward stability, and data-governance evidence. The resulting system avoids two opposite errors: presenting any fitted probability as validated market skill, and making the product unusable because every mathematically coherent model that fails one performance gate is treated as nonexistent. The paper defines the event target, Bayesian reference model, calibration protocol, tiering logic, retraining lifecycle, model freshness, explicit replacement, and bounded Live adaptation. A reference implementation in FinCompass demonstrates four long-horizon Bayesian baseline models, four regime-aware Bayesian research variants, and one stronger validated-research ensemble. The empirical results are reported as software/research evidence, not as proof of durable alpha.
\end{abstract}

\noindent\textbf{Keywords:} Bayesian logistic regression; probabilistic forecasting; calibration; Brier score; locked test; model governance; temporal leakage; online learning; model freshness.

\section{Motivation}
A forecasting application can fail users even while being statistically cautious. Suppose valid historical data are present, a regularized Bayesian probability model can be fitted, and the posterior is numerically stable. If the software nevertheless calls the model ``rejected'' because a bootstrap lower bound, ROC AUC threshold, or walk-forward skill threshold is not met, the user may receive no model and no forecast at all. That behavior confuses \emph{lack of strong demonstrated predictive skill} with \emph{lack of a valid probability model}.

The opposite mistake is also serious: every probability emitted by a fitted model cannot be treated as validated market evidence. A system that returns $0.63$ without a precise target, chronology, calibration, held-out evaluation, and applicability contract merely attaches a percent sign to a score.

The architecture developed here therefore separates three questions:
\begin{enumerate}[itemsep=0.2em]
\item Is the model/data construction valid enough to produce a probability estimate?
\item Has the model demonstrated stronger out-of-sample predictive skill?
\item Does the data provenance support stronger market-grade claims?
\end{enumerate}
This distinction is consistent with the broader literature on proper probability scoring, calibration, temporal validation, data snooping, and prequential evaluation \citep{brier1950verification,gneiting2007proper,gneiting2007calibration,white2000reality,dawid1984prequential}.

\section{Forecast target}
Let $i$ denote an instrument, $b$ a benchmark, $t$ a forecast time, and $H$ a fixed horizon. Define
\begin{align}
R_{i,t}^{(H)} &= \frac{P_i(\tau_H(t))}{P_i(t)}-1,\\
R_{b,t}^{(H)} &= \frac{P_b(\tau_H(t))}{P_b(t)}-1.
\end{align}
For hurdle $\eta$, the target is
\begin{equation}
Y_{i,t}^{(H,\eta)}=\I\left\{R_{i,t}^{(H)}-R_{b,t}^{(H)}>\eta\right\}.
\end{equation}
The desired probability is
\begin{equation}
p_{i,t}=\Prb\left(Y_{i,t}^{(H,\eta)}=1\mid\mathcal{F}_t\right),
\end{equation}
where $\mathcal{F}_t$ contains only information that was actually available at time $t$.

The forward label does not exist until $\tau_H(t)$ exists. This maturity condition is central to both offline retraining and Live adaptation. Re-running software earlier cannot create a label that the market has not yet produced.

\section{Bayesian reference model}
For a standardized feature vector $x$, FinCompass v2.0 uses a regularized Bayesian logistic reference model
\begin{equation}
\Prb(Y=1\mid x,\beta)=\sigmoid(\beta_0+x^\top\beta).
\end{equation}
For slope coefficients,
\begin{equation}
\beta_j\sim\mathcal{N}(0,\tau^2),
\end{equation}
with a wider prior for the intercept. The implementation uses a maximum-a-posteriori estimate and a Laplace approximation to the posterior covariance. This gives an interpretable, strongly regularized model whose uncertainty can be propagated to forecast probabilities without requiring a neural network or a large black-box architecture.

The initial reference feature contract is deliberately compact. For monthly equity forecasting it contains relative momentum over multiple lookbacks, absolute and benchmark volatility, drawdown, distance from a trailing high, a moving-average trend measure, momentum acceleration, downside volatility, benchmark correlation, rolling beta, and trend slope. These features are backward-looking. The broader v2.0 analytics library is not automatically injected into the model.

\section{Calibration}
A logistic model's raw score is not automatically a well-calibrated probability in finite samples. FinCompass therefore fits a separate calibrator using a later validation partition. The locked test remains untouched until model construction is complete. Calibration quality is assessed with expected calibration error (ECE) and calibration slope/intercept in addition to proper scores.

For observed binary outcome $y\in\{0,1\}$ and probability $p$, the Brier score is \citep{brier1950verification}
\begin{equation}
BS=(p-y)^2.
\end{equation}
Log loss is
\begin{equation}
LL=-[y\log p+(1-y)\log(1-p)].
\end{equation}
Both are strictly proper scoring rules when used appropriately \citep{gneiting2007proper}.

\section{Hard validity and evidence tiers}
The key v2.0 change is to separate \emph{hard validity} from \emph{predictive-evidence promotion}.

\subsection{Hard validity}
A Bayesian reference model is unavailable if any indispensable condition fails. The implementation checks, among other conditions, that train/validation/test partitions have adequate rows and both classes, the locked test spans multiple dates, model coefficients and posterior covariance are finite, probabilities are finite and lie in $(0,1)$, chronology is ordered, and test targets are matured forward outcomes.

These are conditions for a coherent model, not claims of market skill.

\subsection{Bayesian baseline}
If hard validity passes but the stronger locked-test gate does not, the model receives the tier
\begin{center}
\texttt{bayesian\_baseline}.
\end{center}
The model may produce a Forecast with an explicit ``Limited evidence'' label. It cannot control the adaptive Live probability.

\subsection{Validated research}
The existing stronger gate remains unchanged. It includes minimum locked-test sample and temporal support, ROC AUC, Brier skill, log-loss skill, ECE, calibration slope, moving-date-block bootstrap bounds, and walk-forward stability. A passing model on real historical data with incomplete historical-universe governance receives \texttt{validated\_research}.

\subsection{Validated market}
A stronger \texttt{validated\_market} tier additionally requires documented point-in-time features, survivorship control, delistings, corporate-action treatment, and supporting evidence. A set of Boolean assertions without evidence is insufficient.

\section{Why a weak baseline can still exist}
Consider a reference model whose test AUC is below $0.5$ and whose Brier skill is slightly negative. It should not be described as having demonstrated predictive advantage. Yet if the model is properly constructed, calibrated as far as the data support, and numerically coherent, the correct conclusion is not necessarily that no probability model exists. The conclusion is that the estimate has \emph{limited demonstrated predictive evidence}. The user can inspect it, while the software clearly prevents stronger claims.

This distinction is particularly useful in markets, where strong and stable out-of-sample skill is rare and repeated model search can itself create false discoveries \citep{white2000reality,harvey2016crosssection,bailey2017pbo}.

\section{Enhanced model and promotion}
The existing FinCompass enhanced model combines a Bayesian logistic component, histogram gradient boosting, and a random forest, with separated component calibration, non-negative stacking, and final calibration. Complexity does not automatically make it superior. The enhanced candidate must earn promotion under the stronger locked-test gate.

The hierarchy is therefore
\begin{equation}
\text{unconditional event rate}\rightarrow\text{Bayesian reference}\rightarrow\text{enhanced model}.
\end{equation}
The enhanced model should be judged against both the unconditional baseline and the simpler Bayesian reference, not merely against chance ranking.

\section{Temporal validation and dependence}
Long-horizon labels overlap. Random row splitting can therefore contaminate training and evaluation. FinCompass uses chronological partitions with forward-target purge and embargo concepts. The locked test is not used for feature selection, calibration choice, ensemble weighting, or threshold tuning. Moving-date-block bootstrap resampling preserves same-date cross-sectional clustering and part of the serial dependence induced by overlapping horizons \citep{kunsch1989jackknife,politis1994stationary}.

\section{Model freshness and delta retraining}
A valid old model is not the same as a current model. FinCompass stores the model creation time, training cutoff, validation cutoff, benchmark cutoff, feature contract, applicability domain, and dataset hash. The data system updates history incrementally by fetching a missing tail plus overlap; overlap revisions are journaled rather than silently replacing history.

Retraining uses the accumulated research corpus,
\begin{equation}
D_{new}=D_{historical}\cup D_{verified\,delta}\cup D_{accepted\,revisions},
\end{equation}
not the delta alone. Only observations whose forward targets have matured may enter the supervised labels.

\section{Candidate replacement}
If model $M_1$ is active and retraining creates $M_2$, then $M_1$ remains active through the states \texttt{not\_ready}, \texttt{failed}, \texttt{rejected}, and \texttt{validated}. A validated candidate is offered to the user; it is not silently activated. This preserves reproducibility and prevents ``newest model wins'' behavior.

A particularly important downgrade rule follows from the new tier hierarchy: an active \texttt{validated\_research} model must not be silently replaced by a newer model that is only \texttt{bayesian\_baseline}. Newer data do not compensate for weaker evidence without explicit review.

\section{Live adaptation}
Live adaptation remains separate from the baseline tier. A \texttt{bayesian\_baseline} forecast may be tracked for freshness, but it does not control adaptive probability updates. Stronger validated models may serve as Live anchors. The adaptive residual is updated only after the exact forward outcome has matured, following a prequential logic \citep{dawid1984prequential}. If freshness, calibration, temporal breadth, loss-regret, or drift conditions fail, the system falls back to the frozen anchor.

\section{Reference implementation results}
A private/reference FinCompass v2.0 build was used to test the architecture on the historical monthly equity sample already used by the project. The source contains eight US equities (IBM, AAPL, MSFT, XRX, AMZN, DELL, GOOGL, ADBE) and the S\&P 500 benchmark. The source ends in June 2022. It does not provide survivorship-complete or market-grade point-in-time universe evidence, so these experiments must not be interpreted as proof of deployable alpha.

\begin{table}[ht]
\centering
\caption{Bayesian reference-family locked-test results. A baseline tier means hard validity passed while the stronger research gate did not.}
\begin{tabular}{rrrrl}
\toprule
Horizon & ROC AUC & Brier skill & Test evidence & Tier\\
\midrule
6 months  & 0.419 & -0.023 & hard-valid & Bayesian baseline\\
12 months & 0.480 & -0.012 & hard-valid & Bayesian baseline\\
24 months & 0.415 & -0.004 & hard-valid & Bayesian baseline\\
36 months & 0.303 & -0.108 & hard-valid & Bayesian baseline\\
\bottomrule
\end{tabular}
\end{table}

The 12-month Bayesian baseline used 532 locked-test observations spanning 76 distinct observation dates and 2,283 calendar days. Its ECE was approximately 0.035. However, its AUC was below $0.5$ and the fitted calibration slope was negative; accordingly, the stronger research gate did not pass. This is exactly the case the new architecture is designed to represent honestly: a mathematically usable reference probability model with limited predictive evidence, not a validated skill claim.

The existing enhanced 12-month model remains stronger evidence within the same limited source context: it previously passed the configured \texttt{validated\_research} gate with ROC AUC approximately 0.576, positive Brier skill approximately 0.036, positive log-loss skill approximately 0.027, and ECE approximately 0.045. The result does not establish universal market validity and remains bounded by the dataset's provenance and universe limitations.

\subsection{Regime-aware Bayesian research variants}
Version 2.0 also evaluates a regime-aware extension rather than assuming that the same feature relationships hold identically in every market environment. A three-state Gaussian hidden Markov model is fitted to backward-looking benchmark features and filtered forward so later observations cannot rewrite earlier regime assignments. The filtered state probabilities are then supplied to the regularized Bayesian reference model as additional covariates. The regime layer is a research candidate, not an automatic promotion.

\begin{table}[ht]
\centering
\caption{Regime-aware Bayesian locked-test results. All four remain baseline-tier research models under the current evidence contract.}
\begin{tabular}{rrrrl}
\toprule
Horizon & ROC AUC & Brier skill & ECE & Tier\\
\midrule
6 months  & 0.444 & -0.045 & 0.086 & Bayesian baseline\\
12 months & 0.508 & -0.059 & 0.067 & Bayesian baseline\\
24 months & 0.485 &  0.009 & 0.029 & Bayesian baseline\\
36 months & 0.321 & -0.043 & 0.163 & Bayesian baseline\\
\bottomrule
\end{tabular}
\end{table}

The 24-month regime variant produced slightly positive Brier and log-loss skill in this locked sample, but its calibration slope remained negative and the overall evidence contract did not justify promotion. This is precisely why the model ladder is comparative: additional complexity is retained for research only when it earns stronger evidence.


\section{Interpretation for users}
The intended product language is deliberately different by tier.

For a Bayesian baseline:
\begin{quote}
Forecast available --- Limited evidence. This probability comes from the standard Bayesian reference model. The model is mathematically usable, but it has not demonstrated enough out-of-sample improvement to qualify for the stronger validated-research tier.
\end{quote}

For a stronger model:
\begin{quote}
Validated research model. This model passed the configured locked-test probability and stability gates for its documented dataset and applicability domain.
\end{quote}

The software should prefer the stronger exact-domain model automatically, but it should not leave the user with nothing merely because only a hard-valid baseline is currently available.

\section{Limitations}
The reference experiments described here use a small current/surviving set of equities and a historical sample with incomplete survivorship, delisting, and corporate-action governance. The sample also ends well before the current market date. Model freshness warnings are therefore mandatory. The results are useful for validating architecture and evidence-tier behavior, not for claiming current trading performance.

The baseline tier also creates a communication burden: users may mistake any numeric probability for evidence of predictive edge. The UI must therefore keep the Limited-evidence label adjacent to the probability and prevent adaptive Live activation from that tier.

\section{Conclusion}
A usable probabilistic forecasting system should neither accept every model nor discard every model that fails to demonstrate strong alpha-like skill. Separating hard validity from predictive-evidence promotion provides a more faithful Bayesian architecture. It allows a regularized, calibrated reference model to exist with limited claims, preserves stricter validation for stronger tiers, and keeps Live control behind stronger evidence. In a noisy nonstationary domain, the strength of the claim should be governed by the strength of the evidence.

\bibliographystyle{plainnat}
\bibliography{refs}
\end{document}
